% Transcription — fonds Alexandre Grothendieck, shelfmark 14, batch 4 (pages 61-80).
% Established from the facsimile archives/batches/14/batch-04.pdf (a mirror of
% https://grothendieck.umontpellier.fr/14.pdf), per the skill
% transcribe-grothendieck. The batch file carries no cover sheet: PDF page N
% is page 60+N of the folder. The archivists' pencil numbers bottom-left
% agree on every leaf (61 on the first, 80 on the last).
%
% Pass: Opus 5.5 (claude-opus-5-5), 2026-09-23 — first pass, unchecked against the pages by a human.
%
% Transcribed: 61-70, 72, 73, 75, 77, 79, 80 (sixteen pages). Skipped: 71
% (blank; 72 shows through), 74 (an English typescript on Hopf algebra /
% coalgebra structures on K[t], page « 7 », scanned upside down), 76 and 78
% (a French typescript on topos cohomology, pages « -20- » and « -25- »,
% Corollaire 2.6 and Définition 3.5 on flasque sheaves; 78 upside down).
% None carries his hand.
%
% What the batch is. Two runs, neither paginated by him, no date.
%   61-70   the end of a numbered sequence on correspondences between X and
%           the Albanese/Picard pair (A_X, B_X): Cor. 2.2 (^t(u_Z) = u_{tZ}),
%           a bracketed Prop. 2.3 « devrait remonter au n° 1 », Prop. 2.4
%           (if a(γ(Z)) : H^1(Y) → H^{2n-1}(X)(n-1) is an isomorphism, i.e.
%           u_Z an isogeny, its inverse is given by a Q-divisor D on X × Y,
%           with denominators controlled by Ker u), Cor. 2.5 (integral
%           version, ℓ prime to N), Cor. 2.6 (hard Lefschetz for H^1 by
%           η = ξ_1…ξ_{n-1}, inverse given by an element of A^1(X × X) ⊗ Q,
%           « Lefschetz vachement » in the margin), the proof via
%           B_X → A_X being a polarisation (restriction to a curve section,
%           equivalent conditions (i)-(iv)), Prop. 2.7 (the converse, D ↦ Z)
%           and its proof via a polarisation ξ (diagram with Ψ_X, λ, μ, φ_X);
%           70: two questions on integral cycles (rational point of the
%           generic fibre of Sym^n(X) → A_X).
%   72-73   loose sheets: H^1(X, μ_N) = _N Pic(X), T_ℓ(Pic^0), the exact
%           sequences for H^{2n-1}(X, Z/N) and H^{2n-1}(X, Z_ℓ)(n-1) with
%           (NS(X)[ℓ])^∨ and H^1(B_X, Z_ℓ); 73 « Relations entre Alb^0_X,
%           Pic^00_X et H^1(X), H^{2n-1}(X) », up to torsion.
%   75-80   a new run numbered « I », with circled (1)(2)(3): an axiomatic of the
%           category VA of abelian varieties up to isogeny (duality functor
%           *, Chow ring C^*, degree, covariance, projection formula, the
%           « application fondamentale » C^1(A) → Hom(A, A^*)^alt); 77 a
%           second (1) (ℓ-adic realisation T_ℓ, H^1, H^* = ΛH^1, compatibility
%           with duality, cycle map γ^{2i}_A) and 2 (η_A : H^{2n}(A) ≃
%           Q_ℓ(-n), « vérifier s'il n'y a pas un signe ! »); 79-80 3)
%           u : A' → A, u_*(1_A) = d·1_{A'}, the class of the image cycle,
%           in terms of T_ℓ. Page 80 breaks off mid-sentence (« … forme
%           2(n'-n)-fois linéaire alternée sur T_{Q_ℓ}(A') »).
%
% Notation fixed for the batch. A_X his Albanese (A_X = B_X^*), B_X his
% Picard variety (B_X = Pic^00_{X/k}); u_Z : B_X → A_Y the homomorphism
% defined by a correspondence Z (his 2.1, previous batch); a(γ(Z)) the
% cohomological effect of the class of Z; A^1, A_1 his groups of cycles
% (codim. 1, dim. 1); \mathcal{C}\ell his « Cl » (divisor classes);
% \mathcal{VA} his script VA; \mathcal{C}^* his rounded C (Chow ring).
%
% The dating is the inventory's own for the folder, copied verbatim. Its
% group, [10-18] « Théorie des motifs », is dated
% [à partir de 1958-à partir de 1982].
\documentclass[11pt,a4paper]{article}
% Le préambule commun à toutes les transcriptions du fonds Grothendieck.
%
% Il tient en une idée : **l'incertitude se compose, elle ne se résout pas.**
% Une transcription qui lisse un mot douteux a détruit la seule chose pour
% laquelle on la faisait — un lecteur doit pouvoir savoir, à chaque endroit,
% ce qui était sur le feuillet et ce qui a été ajouté. Les six macros qui
% suivent sont donc l'appareil critique, et non de la décoration.
%
% Les mêmes commandes sont reconnues par `scripts/render.mjs`, qui produit la
% vue de lecture du site. Une macro ajoutée ici doit l'être aussi là-bas,
% faute de quoi le rendu échouera bruyamment — ce qui est voulu : un rendu
% silencieusement faux est pire qu'un rendu absent.

\ProvidesPackage{grothendieck}[2026/08/08 Transcriptions du fonds Grothendieck]

\usepackage{amsmath,amssymb,amsthm}
\usepackage{mathrsfs}
\usepackage{stmaryrd}
% Les diagrammes commutatifs sont partout dans ce fonds : c'est la langue
% même de la géométrie algébrique de Grothendieck, et une transcription qui
% les rendrait en prose ne transcrirait rien.
\usepackage{tikz-cd}
% Français par défaut — c'est la langue des feuillets — mais l'anglais est
% chargé aussi : la traduction et le résumé sont des documents anglais, et la
% typographie française (espaces avant les deux-points) y serait une faute.
% Ces documents basculent par \selectlanguage{english} en tête de corps.
\usepackage[english,french]{babel}
\usepackage{fontspec}
\usepackage{geometry}
\geometry{a4paper,margin=25mm,marginparwidth=28mm}
\usepackage{marginnote}
\usepackage{xcolor}
% ulem, not soul. `soul`'s \st and \ul are built on a character-by-character
% scan that XeLaTeX's font machinery breaks: the document fails with an
% undefined control sequence rather than degrading. ulem does the same two jobs
% and is Unicode-engine safe. `normalem` stops it hijacking \emph, which the
% transcriptions use for Grothendieck's own emphasis.
\usepackage[normalem]{ulem}
\usepackage{hyperref}
\hypersetup{colorlinks=true,linkcolor=black,urlcolor=black,pdfborder={0 0 0}}

% --- Métadonnées du lot -----------------------------------------------------
% Reprises telles quelles par le script de rendu, qui les affiche en tête de
% la vue de lecture. Elles fixent ce que le fichier prétend couvrir : sans
% elles, une transcription flotte sans point d'attache dans un fonds de seize
% mille feuillets.

\newcommand{\folder}[1]{\gdef\grothFolder{#1}}
\newcommand{\batch}[1]{\gdef\grothBatch{#1}}
\newcommand{\leaves}[2]{\gdef\grothFirst{#1}\gdef\grothLast{#2}}
\newcommand{\foldertitle}[1]{\gdef\grothFoldertitle{#1}}
\gdef\grothFolder{?}\gdef\grothBatch{?}\gdef\grothFirst{?}\gdef\grothLast{?}\gdef\grothFoldertitle{}

% --- Le résumé --------------------------------------------------------------
%
% Il ouvre la lecture modernisée, sous le titre « Résumé », et il est composé
% en petit. Ce n'est pas un ornement : c'est le seul passage du document qui
% ne soit pas des mathématiques, et un lecteur doit pouvoir le reconnaître
% avant de l'avoir lu — puis le sauter s'il connaît déjà le sujet, ou s'y
% arrêter s'il ne le connaît pas. Le filet qui le suit marque la reprise du
% corps.

\newenvironment{resume}
  {\small\setlength{\parskip}{0.45em}}
  {\par\medskip\hrule\medskip}

% --- Mots-clés ---------------------------------------------------------------
%
% En anglais, en clôture du résumé de la lecture modernisée : le vocabulaire
% moderne sous lequel chercher ce que les feuillets construisent. Le manifeste
% du site (`npm run manifest`) les extrait pour étiqueter la cote entière —
% cette ligne est la source unique des tags, il n'y a pas de fichier à côté.
\newcommand{\keywords}[1]{\par\smallskip\noindent
  {\footnotesize\textsc{Keywords} --- \textit{#1}}\par}

% --- L'appareil critique ----------------------------------------------------

% Le numéro de feuillet, en marge. C'est l'ancre qui permet de régler un
% désaccord sur une formule en regardant *un* feuillet plutôt qu'une chemise
% de six cents. Le site s'en sert aussi pour tourner les pages du fac-similé
% quand on fait défiler la transcription.
\newcommand{\leaf}[1]{%
  \marginnote{\footnotesize\color{gray!70}\textsf{#1}}%
  \label{leaf:#1}\ignorespaces}

% Illisible : on ne devine pas. Un mot inventé qui se lit comme les autres est
% le pire résultat possible.
\newcommand{\ill}{\textcolor{red!60!black}{[\,\ldots\,]}}

% Lecture douteuse : le mot est donné, et signalé comme incertain.
\newcommand{\uncertain}[1]{\uline{#1}}

% Ajout de l'éditeur — une lettre manquante, un mot sous-entendu.
\newcommand{\add}[1]{\textcolor{blue!50!black}{[#1]}}

% Biffé par Grothendieck lui-même. Ce qu'il a rayé fait partie du document :
% une rature dit souvent où la pensée a changé de direction.
\newcommand{\struck}[1]{\sout{#1}}

% Note du transcripteur, sur l'état matériel du feuillet.
\newcommand{\note}[1]{\par\smallskip{\small\color{gray!60!black}\textit{[#1]}}\par\smallskip}

% Note marginale de Grothendieck, distincte du corps du texte.
\newcommand{\marginal}[1]{\par\smallskip\noindent\hspace{1em}%
  {\small\color{gray!60!black}#1}\par\smallskip}

% --- Titre ------------------------------------------------------------------

\AtBeginDocument{%
  \begin{center}
    {\large\bfseries Fonds Alexandre Grothendieck --- cote n\textsuperscript{o}~\grothFolder}\\[2pt]
    {\itshape\grothFoldertitle}\\[4pt]
    {\small feuillets \grothFirst--\grothLast\ (lot \grothBatch)}\\[2pt]
    {\footnotesize\color{gray!60!black}Transcription établie d'après le fac-similé numérisé
     par l'Université de Montpellier.\\
     Les lectures incertaines sont soulignées, les illisibles notées [\,\ldots\,],
     les ajouts entre crochets.}
  \end{center}
  \vspace{1em}}

\endinput


\folder{14}
\batch{4}
\pages{61}{80}
\dating{1965-[vers 1971]}
\watermark{Édition de démonstration}
\foldertitle{Théorie des cycles algébriques : conjectures de Hodge, Tate, Lefschetz, Weil : notes manuscrites (s.d.), lettre (1965).}

\begin{document}

\page{61}
\note{la page 60 (lot précédent) s'achève sur un énoncé « Cor. Main 2.2 », barré de grands traits obliques, sur les homomorphismes $B_X \to A_Y$ définis par un cycle sur $X \times Y$ ; le Corollaire 2.2 ci-dessous est un autre énoncé sous le même numéro.}

\textbf{Corollaire 2.2.} Notations de 2.1. Considérons ${}^{t}(u_{Z}) : B_{Y} \to A_{X}$, alors
${}^{t}(u_{Z}) = u_{({}^{t}Z)}$, i.e. \ill{} pour tout diviseur $D$ alg. équiv. à $0$ sur $Y$, tel
que $Z(D)$ soit défini,
\[
  S_{X}(Z(D)) = {}^{t}(u_{Z})(D)
\]
Vérification par exemple par voie cohomologique.
\marginal{Attention aux signes éventuels !}

\marginal{devrait remonter au n° 1.}
[\,\textbf{Proposition 2.3.} Conditions de 1.1. \struck{Alors} Considérons
$\sigma_{Z} = {}^{t}(u_{Z}) : B_{X} \to B_{Y}$. Alors \ill{}, pour tout diviseur alg. équiv. à $0$
$D$ sur $X$, tel que $Z(D)$ soit défini,
\[
  {}^{t}Z(D) = {}^{t}(u_{Z})(D)
\]

N.B. Ces deux propositions servent à une autre démonstration de 1.1, via un homom.
\uncertain{de} schémas en groupes
\[
  \underline{\mathrm{Pic}}_{X/k} \longrightarrow \underline{\mathrm{Pic}}_{Y/k} \ !
\]

Vérification de 2.3 par \uncertain{AQT}.\,]
\marginal{attention aux signes éventuels !}
\note{la note marginale « Attention aux signes éventuels ! » est écrite deux fois, en face du
Corollaire 2.2 et en face de la vérification de 2.3 ; les crochets qui encadrent 2.3 et le N.B.
sont de sa main.}

\textbf{Proposition 2.4.} Avec les notations de 2.1. \struck{Supposons que} Fixons $\ell$,
pour \struck{un $\ell$ convenable} \uncertain{que} l'homom.
\[
  a(\gamma(Z)) : H^{1}(Y, \mathbb{Q}_{\ell}) \longrightarrow H^{2n-1}(X, \mathbb{Q}_{\ell})(n-1)
\]
soit un isom \struck{\ill{}}, \uncertain{i.e.} que $u_{Z}$ soit une isogénie (cond. néc.
\struck{et suffisante}). Alors il existe un \struck{cycle} diviseur

\page{62}
\struck{algébrique} $D$ à coefficients rationnels, \struck{\ill{}} sur $X \times Y$ tel que
l'homom. $a({}^{t}\gamma(D)) : H^{2n-1}(X, \mathbb{Q}_{\ell})(n-1) \to H^{1}(Y, \mathbb{Q}_{\ell})$
soit inverse de $a(\gamma(Z))$, ce qui caractérise $D$ comme élément de
$\mathcal{C}\ell(X \times Y) \otimes \mathbb{Q}$.
L'homomorphisme de $\mathrm{Hom}_{\mathbb{Q}}(A_{Y}, B_{X})$ défini par $D$ n'est autre que
$u_{Z}^{-1}$, \struck{\ill{}}. Pour qu'on puisse choisir $D$ à coefficients entiers, il faut et il
suffit que $u$ soit un isom. \struck{Pour que \ill{}} Plus généralement, si $\operatorname{Ker} u$
est annulé par $N$, \struck{\ill{}} (attention, $\operatorname{Ker} u$ \uncertain{en tant que}
schéma en gr. \uncertain{non réd.}) alors on peut choisir $D$ de la forme $D = N^{-r} D_{0}$,
$D_{0}$ étant un div. à coeff. entiers, et \ill{} de la forme $N^{-r} D_{0}$, \ill{} $r \geqslant 0$,
\struck{\ill{}} \ill{} diviseurs premiers de $n$.
\note{les deux lignes qui suivent « $D_0$ étant un div. à coeff. entiers » sont surchargées
d'ajouts interlinéaires et d'une rature ; seuls les fragments donnés se lisent.}
On écrit le diagramme \uncertain{commutatif}

\begin{tikzcd}
  H^{1}(Y, \mathbb{Q}_{\ell}) \arrow[r, "a(\gamma(Z))"] & H^{2n-1}(X, \mathbb{Q}_{\ell})(n-1) \arrow[d, "\wr"] \\
  H^{1}(A_{Y}, \mathbb{Q}_{\ell}) \arrow[u, "\wr"] \arrow[r, "u^{(1)}"] & H^{1}(B_{X}, \mathbb{Q}_{\ell})
\end{tikzcd}

où toutes les flèches \uncertain{écrites} sont $\mathbb{Z}$-algébriques, les flèches verticales
des isom. Ceci prouve que

\page{63}
\marginal{(\uncertain{$b$} étant donné par $\operatorname{card} \operatorname{Ker} u$)}
\note{ajout encadré en tête de page ; il se rapporte au « meilleur dénominateur possible $b$ » quelques lignes plus bas.}

$a(\gamma(Z))$ soit un isom., il faut et il suffit que $u^{(1)}$ le soit, i.e. que $u$ soit une
isogénie. Alors $u$ admet un inverse unique dans $\mathrm{Hom}_{\mathbb{Q}}(A_{Y}, B_{X})$,
\struck{défini} dont l'homom. inverse de $u^{(1)}$ est défini par un \ill{}
$\mathbb{Q}$-cycles algébriques. \struck{à coeff. en} $R$, le meilleur dénominateur possible
$b$ ($D$ \ill{} $a(\gamma(Z))^{-1}$) est défini par un $\mathbb{Q}$-cycle algébrique, ce qui
prouve la première assertion, \struck{La seconde} \uncertain{ayant} \struck{\ill{}}
\struck{\ill{}} dénominateur que $R$. Cela prouve la première assertion, la seconde résulte de
la connaissance de l'effet cohomologique d'une donnée de corr. \struck{de} diviseurs
\uncertain{associée}, via l'effet sur les VA … (cf …). Il reste :
\struck{prouver} \uncertain{notons} que si $D$ admet un certain dénominateur,
\uncertain{inverse}, un dénominateur valant pour le $\mathbb{Q}$-hom $A_{Y} \to B_{X}$ inverse de
$u$, \ill{} il \struck{\ill{}} un \ill{} par \ill{} meilleur \uncertain{puissance} de
\struck{cardt} $N$.
\note{les points de suspension « sur les VA … (cf … ) » sont de lui : la référence est laissée
en blanc. La fin du paragraphe, à partir de « Il reste », est d'une écriture rapide et raturée.}

\textbf{Corollaire 2.5.} Avec les notations de 2.4, \uncertain{pour}

\page{64}
\uncertain{que} $u$ soit une isogénie et que $N$ soit premier à $\ell$ (où $\ell \neq$ car $k$)
il f. et s. que le composé
\[
  H^{1}(Y, \mathbb{Z}_{\ell}) \longrightarrow H^{2n-1}(X, \mathbb{Z}_{\ell})(n-1) \longrightarrow
  \underline{H^{2n-1}(X, \mathbb{Z}_{\ell})(n-1)}
\]
soit \struck{surjectif} bijectif ; \struck{et (il suffit \ill{}} \struck{(Pour que de plus
$H^{2n-1}(X, \mathbb{Z}_{\ell})$ \uncertain{soit sans torsion} \ill{})} que
$H^{1}(Y, \mathbb{Z}_{\ell}) \to H^{2n-1}(X, \mathbb{Z}_{\ell})(n-1)$ soit bijectif.
\note{« u soit une isogénie et que » est ajouté au-dessus de la première ligne et rattaché par
un trait au début de l'énoncé, qui commence à la dernière ligne de la page 63. Le second
$H^{2n-1}$ du composé est souligné par lui : le soulignement ne se lit pas avec certitude
(quotient par la torsion ?). La parenthèse interlinéaire sur la torsion paraît traversée par le
même trait que « et (il suffit » ; l'ordre exact des ratures est incertain.}

\marginal{\emph{Th} ? C'est un cas particulier important du « Lefschetz
\uncertain{vachement} » donnant « Lefschetz » aux \uncertain{surfaces} …}

\textbf{Corollaire 2.6.} Soit $X$ projective lisse conn. de dim $n \geqslant 2$, et
$\xi_{1}, \ldots, \xi_{n-1} \in$ \struck{$\mathrm{Pic}(X)$ \ill{} des \ill{}}
[$\xi_{1}, \ldots, \xi_{n-1} \in A^{1}(X)$ leurs classes.] \struck{provenant} des faisceaux
amples, Alors pour tout entier $\ell$, l'hom. cup-produit par
$\eta = \xi_{1}\xi_{2}\cdots\xi_{n-1}$ :
\[
  H^{1}(X, \mathbb{Q}_{\ell}) \longrightarrow H^{2n-1}(X, \mathbb{Q}_{\ell})(n-1)
\]
est un \emph{isomorphisme}. L'isom. inverse est définissable par un élément de
\struck{\ill{}} $A^{1}(X \times X) \otimes \mathbb{Q}$, \struck{i.e. de} ou encore de
$\mathcal{C}\ell(X, X) \otimes \mathbb{Q}$, et \struck{cet élément} cet élément est unique et ne
dépend pas de $\ell$.
\note{la mention encadrée « $\xi_1, \ldots, \xi_{n-1} \in A^1(X)$ leurs classes » est écrite au-dessus
de la ligne biffée ; « Alors pour tout » est un ajout encadré.}

Il faut prouver que l'hom de VA
\[
  u : B_{X} \longrightarrow A_{X}
\]

\page{65}
« correspondant » est une isogénie, ou, \uncertain{mieux}, a un noyau fini. Pour bien faire,
\struck{il faudrait} on peut prouver mieux, savoir qu'il correspond à une
\emph{polarisation} de $B_{X}$ (compte tenu de $A_{X} \simeq B_{X}^{*}$).

\struck{On est \ill{} \ill{}, peut \ill{} le \ill{} de \ill{} \ill{} comme composé d'hom. de
restriction}

\struck{$B_{X} = B_{X_{0}} \xrightarrow{\rho_{0}} B_{X_{1}} \xrightarrow{\rho_{1}} B_{X_{2}} \cdots$}
\struck{$\xrightarrow{\rho_{n-3}} B_{X_{n-2}} \xrightarrow{\sigma} A_{X_{n-2}} \xrightarrow{i} \cdots$} \struck{\ill{}}

\[
  B_{X} \xrightarrow{\ \rho\ } B_{Y} \xrightarrow{\ \lambda\ } A_{Y} \xrightarrow{\ \sigma\ } A_{X}
\]
$Y$ \uncertain{étant} défini t. comme l'intersection de $n-2$ \struck{hyperplans} diviseurs
« génériques » \uncertain{convenables} \uncertain{correspondant} : $L_{1}, \ldots, L_{n-2}$
(supposons \emph{très} amples), $\rho$ est la restriction, $\sigma$ l'injection, et $\lambda$ est
défini comme $u$, \uncertain{avec} $L_{n-1}|Y$ sur $Y$. On \uncertain{sait} \ill{} \ill{}
(\uncertain{des propriétés des} \uncertain{pinceaux} …) que $\rho$, $\sigma$ sont des isogénies,
\struck{et} \struck{(d'ailleurs transposées l'une de l'autre)},
\note{tout le passage depuis « On est » jusqu'à « l'une de l'autre) » est annulé par deux
grands traits en croix et un trait vertical dans la marge gauche ; la ligne-diagramme
$B_X = B_{X_0} \to \cdots$ est en outre barrée d'un trait horizontal et se termine par une
surcharge illisible. Seul le diagramme $B_X \to B_Y \to A_Y \to A_X$ reste hors des ratures
propres à la ligne.}

Or une fonctorialité évidente en $X$, $Y$ permet d'interpréter de la façon suivante l'hom.
$B_{X} \to A_{Y}$ défini dans 2.1, lorsque $Z$ est un morphisme
\uncertain{défini} par un

\page{66}
\[
  v : C \longrightarrow X \times Y ,
\]
où $C$ est une courbe lisse projective, donnée par
\[
  v_{1} : C \longrightarrow X , \qquad v_{2} : C \longrightarrow Y .
\]
Il suffit dans ce cas de prendre le composé
\[
  B_{X} \xrightarrow{\ v_{1}^{*}\ } B_{C} \xrightarrow{\ \text{dualité}\ } A_{C}
  \xrightarrow{\ v_{2*}\ } A_{Y} .
\]
Si on exprime en termes de données de correspondances divisorielles sur $B_{X} \times B_{Y}$, il
suffit de prendre la donnée \struck{\uncertain{dualisante}} canonique sur $B_{C} \times B_{C}$,
avec les hom. $B_{X} \xrightarrow{v_{1}^{*}} B_{C}$, $B_{Y} \xrightarrow{v_{2}^{*}} B_{C}$.
Ceci posé, si $X = Y$, $v_{1} = v_{2}$, on voit que \struck{cela donc \ill{}} les conditions
suivantes sont équivalentes :
\begin{enumerate}
  \item[(i)] L'hom. $B_{X} \to A_{X} = B_{X}^{*}$ est un hom. de polarisation.
  \item[(ii)] L'hom. $B_{X} \to B_{X}^{*} = A_{X}$ est une isogénie.
  \item[(iii)] L'hom. $v^{*} : B_{X} \to B_{C}$ est injectif (mod gpes finis).
  \item[(iv)] L'hom. $v_{*} : A_{C} \to A_{X}$ est surjectif.
\end{enumerate}
\note{« (mod gpes finis) » est placé entre (iii) et (iv), dans une accolade qui paraît valoir
pour les deux.}

\page{67}
D'ailleurs si $\xi \in A_{1}(X)$ est défini par $u : C \to X$, alors $B_{X} \to A_{X}$ est aussi,
du point de vue cohomologique, défini par le \uncertain{cup-produit} par
\[
  \gamma(c) : H^{1}(A_{X}, \mathbb{Q}_{\ell}) \simeq H^{1}(X, \mathbb{Q}_{\ell})
  \longrightarrow H^{1}(B_{X}, \mathbb{Q}_{\ell}) \xleftarrow{\ \sim\ }
  H^{2n-1}(X, \mathbb{Q}_{\ell})(n-1) .
\]
On est ramené au fait bien connu que si $C$ est une courbe sur $X$, intersection « gén. » de
diviseurs très amples, alors $A_{C} \to A_{X}$ est surjectif.
\note{cette dernière phrase est séparée de ce qui précède par un trait oblique et encadrée.}

\textbf{Proposition 2.7.} $X$ et $Y$ sont comme d'habitude, $D \in A^{1}(X \times Y)$, d'où
\[
  u_{D} : A_{Y} \longrightarrow B_{X}
\]
et
\[
  a(\gamma(D)) : H^{2n-1}(X, \mathbb{Q}_{\ell})(n-1) \longrightarrow H^{1}(Y, \mathbb{Q}_{\ell}) .
\]
Pour que $a(\gamma(D))$ soit un isom., il f. et s. que $u_{D}$ soit une isogénie. Alors il existe
$Z \in A_{1}(X \times Y) \otimes \mathbb{Q}$ tel que
\[
  a(\gamma(Z)) : H^{1}(Y, \mathbb{Q}_{\ell}) \longrightarrow
  H^{2n-1}(X, \mathbb{Q}_{\ell})(n-1)
\]
soit \struck{défini} inverse de $a(\gamma(D))$ (quel que soit $\ell$). \struck{Ce $Z$ est}
\note{devant le but $B_X$ de $u_D$, un premier symbole est raturé ; devant $a(\gamma(Z))$, un petit « t » en exposant est biffé.}

\page{68}
\struck{Rappelons le $\mathbb{Q}$-hom} \struck{On écrit le diagramme \ill{}}
\[
  u_{Z} : B_{X} \longrightarrow A_{Y} \quad (\text{une isog.})
\]
défini par $Z$ au moyen de 2.1, n'est autre que $u_{D}^{-1}$.

\emph{Dém.} On utilise le diagramme commutatif

\begin{tikzcd}
  H^{2n-1}(X, \mathbb{Q}_{\ell})(n-1) \arrow[r, "a(\gamma(D))"] \arrow[d, "\Psi_{X}"'] & H^{1}(Y, \mathbb{Q}_{\ell}) \\
  H^{1}(B_{X}, \mathbb{Q}_{\ell}) \arrow[r, "u_{D}^{(1)}"] & H^{1}(A_{Y}, \mathbb{Q}_{\ell}) \arrow[u, "\Phi_{Y}"']
\end{tikzcd}

\note{les deux flèches verticales portent en outre le signe $\wr$ (isomorphismes) ; à côté de
$\Psi_X$, un mot est biffé.}

Ceci prouve que $a(\gamma(D))$ est un isom. ssi $u_{D}^{(1)}$ l'est, i.e. $u_{D}$ est une
isogénie. Pour conclure à l'existence de $Z$, il faut montrer que $\Phi_{Y}^{-1}$ et
$\Psi_{X}^{-1}$ sont $\mathbb{Q}$-algébriques (indépendamment de $\ell$). \struck{R\ill{}}
\struck{le \ill{} défini} \struck{Ceci, \ill{} faut} Pour $\Psi_{X}$, on est donc ramené :
prouver 2.7 dans le cas particulier de l'hom. can.
\[
  H^{2n-1}(X, \mathbb{Q}_{\ell})(n-1) \longrightarrow H^{1}(B_{X}, \mathbb{Q}_{\ell})
\]
défini par $D \in \mathcal{C}\ell(X, B_{X})$, \uncertain{pour trouver} un élément can.
$Z \in A_{1}(X \times B_{X})$.

\page{69}
Or choisissons un $\xi \in A^{1}(X)$ de polarisation. Alors on a un diagramme commutatif

\begin{tikzcd}
  H^{2n-1}(X, \mathbb{Q}_{\ell})(n-1) \arrow[r, "\Psi_{X}"] & H^{1}(B_{X}, \mathbb{Q}_{\ell}) \\
  H^{1}(X, \mathbb{Q}_{\ell}) \arrow[u, "\lambda"] & H^{1}(A_{X}, \mathbb{Q}_{\ell}) \arrow[u, "\mu"'] \arrow[l, leftrightarrow, "\varphi_{X}^{(1)}"]
\end{tikzcd}

\note{toutes les flèches portent le signe d'isomorphisme ($\sim$ ou $\wr$) ; la flèche
verticale de gauche est annotée « cup $\xi^{n-1}$ ». La flèche horizontale du bas a une pointe à
chaque bout, la droite surchargée : le sens n'est pas tranché sur la page.}

où toutes les flèches \struck{sont \ill{}} autres \uncertain{sont} définies par des classes de
cycles alg. (indép. de $\ell$). ($\mu$ étant défini comme \uncertain{$\Psi_{X}\lambda\varphi_{X}$}
(composé) et étant donc un isom. en vertu de \ill{} et de 2.6.)
Donc $\Psi_{X}^{-1} = \lambda\varphi\mu^{-1}$, et tout revient à voir que $\mu^{-1}$ est
algébrique, ce qui provient du fait que $\mu$ est défini par une $\mathbb{Q}$-isogénie
$B_{X} \to A_{X}$ (cf \ill{}), dont on peut prendre l'inverse …

Reste : traiter le cas de $\Phi_{Y}^{-1}$ (où on peut remplacer $X$ par $Y$). On utilise le
\uncertain{même} diagramme, qui donne
\[
  \varphi_{X}^{-1} = \mu^{-1}\Psi_{X}\lambda
\]
donc on est ramené à inverser \uncertain{les} « algébriques » ; or cela est possible grâce au fait
que $\mu$ est associé à une isogénie de VA.
\note{« (cf » est suivi d'un blanc ; les points de suspension après « l'inverse » sont de lui.
Le $\varphi_X$ de cette page correspond au $\Phi_Y$ de la page 68, $Y$ remplacé par $X$.}

\page{70}
\emph{Questions.} 1) Y a-t-il toujours un cycle \emph{entier} qui donne l'inverse pour
$H^{1}(A_{X}) \to H^{1}(X)$ ? Il suffirait qu'on sache que pour $n$ grand convenable
(\uncertain{donc} $n \geqslant d$) la fibre générique de $\mathrm{Sym}^{n}(X) \to A_{X}$ a un point
rationnel.

2) Y a-t-il toujours un cycle \emph{entier} qui donne l'inverse pour
$H^{2n-1}(X)(n-1) \to H^{1}(B_{X})$ ? Si la réponse à 1) était affirmative, il suffirait de
connaître la réponse à 2) dans le cas des VA.

\page{72}
\[
  H^{1}(X, \mu_{N}) \simeq {}_{N}\mathrm{Pic}(X)
\]
\[
  0 \to {}_{N}\mathrm{Pic}^{0}(X) \to H^{1}(X, \mu_{N}) \longrightarrow {}_{N}\mathrm{NS}(X) \longrightarrow 0
\]
\marginal{$X$ propre sur corps $k$ sép. clos, $N$, $\ell$ premiers à car.\ de $k$}
\[
  H^{1}(X, \mathbb{Z}_{\ell}(1)) \xleftarrow{\ \simeq\ }
  T_{\ell}(\underline{\mathrm{Pic}}^{0}_{X/k})
\]
\note{le but s'écrivait d'abord \struck{$H^{1}(X, A_{X})$}, biffé et remplacé par
$T_\ell(\underline{\mathrm{Pic}}^0_{X/k})$ ; un signe raturé suit.}
\marginal{Utiliser le fait que $\mathrm{Pic}^{\tau}(X)/\mathrm{Pic}^{0}(X)$ est \emph{fini}}
\note{dans la seconde note marginale, un premier quotient, avec $\underline{\mathrm{Pic}}_{X/k}$,
est raturé. La première suite commence par un « $0 \to {}_N\mathrm{Pic}^0(X) \to$ » ajouté dans la
marge gauche, d'une autre encre.}

En particulier, si $\underline{\mathrm{Pic}}^{0}_{X/k}$ propre, p. ex. $X$ normal, on trouve,
introduisant $B_{X} = \underline{\mathrm{Pic}}^{00}_{X/k}$ et $A_{X} = B_{X}^{*}$, un
\[
  H^{1}(X, \mathbb{Z}_{\ell}) \xleftarrow{\ \varphi_{X}^{*}\ } H^{1}(A_{X}, \mathbb{Z}_{\ell}) .
\]

\struck{$H^{2n-1}$} Soit maintenant $X$ lisse, purement $n$-dim. Alors
\[
  H^{2n-1}(X, \mathbb{Z}/N\mathbb{Z}) \simeq H^{1}(X, \mathbb{Z}/N\mathbb{Z})^{\vee}(-n)
\]
d'où
\[
  0 \to ({}_{N}\mathrm{NS}(X))^{\vee} \to H^{2n-1}(X, \mathbb{Z}/N\mathbb{Z})
  \to {}_{N}A_{X}(k)(-n) \to 0
\]
avec ${}_{N}A_{X}(k)(-n) \simeq H^{1}(B_{X}, \mathbb{Z}/N\mathbb{Z})(1-n)$.
Passant à la limite sur $\ell^{\nu}$ :
\[
  0 \to (\mathrm{NS}(X)[\ell])^{\vee} \to H^{2n-1}(X, \mathbb{Z}_{\ell})(n-1)
  \to H^{1}(B_{X}, \mathbb{Z}_{\ell}) \to 0
\]
où $\mathrm{NS}(X)[\ell]$ est la comp.\ $\ell$-primaire, et la dernière flèche un hom.\ algébrique.
\note{les exposants de Tate de ces deux suites ont été corrigés sur la page : $(-n)$ biffé après
$({}_N\mathrm{NS}(X))^\vee$, $(1-n)$ surchargeant une première valeur après
$H^1(B_X, \mathbb{Z}/N\mathbb{Z})$, et $(n-1)$ porté en travers de la flèche du milieu de la
dernière suite ; après $(\mathrm{NS}(X)[\ell])^\vee$ et après $H^1(B_X, \mathbb{Z}_\ell)$, un
exposant (dans le second cas $(1-n)$) est raturé.}

\page{73}
\textbf{Relations entre $\underline{\mathrm{Alb}}^{0}_{X}$, $\underline{\mathrm{Pic}}^{00}_{X}$ et
$H^{1}(X)$, $H^{2n-1}(X)$} ($X$ projective lisse conn. de dim $n$ sur $k$ alg. clos).

\begin{align*}
  &\pi_{1}^{\ell}(\underline{\mathrm{Alb}}^{0}_{X}) \xleftarrow[\ \sim\ ]{(t)} \pi_{1}^{\ell}(X)^{\mathrm{ab}}(\ell)
  && \text{i.e. \emph{torsion près}} \quad [\text{c'est toujours surjectif ?}] \\
  &\pi_{1}(\underline{\mathrm{Pic}}^{00}_{X}) \simeq H^{1}(\underline{\mathrm{Alb}}^{0}_{X})(1)
  \xrightarrow[\ \sim\ ]{t} H^{1}(X)(1)
  && \text{i.e. \emph{torsion près}} \quad [\text{c'est toujours injectif}] \\
  &H^{1}(\underline{\mathrm{Alb}}^{0}_{X}) \xrightarrow[\ \sim\ ]{t} H^{1}(X)
  && \text{i.e. \emph{torsion près}} \quad [\text{c'est toujours injectif}] \\
  &H^{1}(\underline{\mathrm{Pic}}^{00}_{X}) \simeq \pi_{1}(\underline{\mathrm{Alb}}_{X})(-1)
  \xleftarrow[\ \sim\ ]{(t)} H^{2n-1}(X)(n-1)
  && \text{toujours}
\end{align*}
\note{au-dessus du premier but $\pi_1^\ell(X)^{\mathrm{ab}}(\ell)$, un ajout rattaché par une
accolade : « $(t)\ H^{2n-1}(X)(n)$ », qui paraît remplacer un symbole raturé après $(\ell)$.
Les points d'exclamation ou d'interrogation qui closent les deux crochets « injectif » sont
surchargés ; devant « toujours », à la dernière ligne, un mot raturé et illisible (\struck{\ill{}}).}

Si on repère les VA par les $H^{1}$,
\[
  \left\{
  \begin{array}{lll}
    \underline{\mathrm{Alb}}_{X} & \text{correspond :} & H^{1}(X) \\
    \underline{\mathrm{Pic}}_{X} & \text{------------} & H^{2n-1}(X)(n-1)
  \end{array}
  \right.
\]

\page{75}
\note{commence ici une nouvelle suite, numérotée par lui « I » (cerclé), avec des alinéas
aux numéros cerclés, rendus ici (1), (2), (3) : une axiomatique de la catégorie des variétés
abéliennes et de leurs cycles.}

\textbf{I.} (1) $\mathcal{VA}$ catégorie (additive) définie \struck{\ill{}} \uncertain{à isog.
près} \struck{\ill{}} sur un corps alg. clos $k$.
\[
  \mathcal{VA}^{\circ} \xrightarrow{\ *\ } \mathcal{VA} \qquad A \rightsquigarrow A^{*}
\]
foncteur « dualité » \struck{\ill{}}
\[
  \mathcal{VA}^{\circ} \xrightarrow{\ C^{*}\ }
  \begin{pmatrix}
    \text{Anneaux gradués comm.,} \\
    \text{à degrés} \geqslant 0, \\
    C^{0} \simeq \mathbb{Z}
  \end{pmatrix}
  \qquad A \rightsquigarrow C^{*} \quad \text{cycles mod équiv. alg.}
\]
Fonction dimension, additive, invariante par dualité.

(2) $C^{n}(A) \xrightarrow{\ \deg\ } \mathbb{Z}$ isom. can., mais pas fonctoriel
(\uncertain{fonct.} pour isom.) ; $C^{0}(A) \xleftarrow{\ \sim\ } \mathbb{Z}$ isom. fonctoriel.
\note{l'étiquette de la première flèche est surchargée ; on lit au-dessous une sorte de
$\int_x$ ou « deg ».}

Loi \emph{covariante} pour les cycles, formule de projection (contenant la formule de
transposition).

On peut en tirer formellement la loi de \uncertain{Pythagore} pour les cycles …

(3) Application fondamentale

\begin{tikzcd}
  C^{1}(A) \arrow[rr] \arrow[dr, dashed] & & \mathrm{Hom}(A, A^{*})^{\mathrm{alt}} \\
  & C^{1}(A \times A) \arrow[ur, dashed] &
\end{tikzcd}

\[
  C^{1}(A \times B) \longrightarrow \mathrm{Hom}(A, B^{*})
\]
(isom. $C^{1}(A \times B)/C^{1}(A) \times C^{1}(B) \simeq \mathrm{Hom}(A, B^{*})$)
\uncertain{on a un choix}

\marginal{N.B. On pourrait définir axiomatiquement $A^{*}$ par cette formule}
\note{les deux flèches obliques du diagramme sont en pointillé sur la page ; le
$C^1(A \times A)$ surcharge une première lettre. Le N.B. est encadré dans le coin inférieur
gauche.}

\page{77}
\note{la page reprend la numérotation de la page 75 : un nouvel alinéa (1), cerclé (réalisation
$\ell$-adique), un second groupe marqué en marge d'un signe illisible, puis un alinéa 2, non
cerclé.}

(1) \[
  \mathcal{VA} \xrightarrow{\ T_{\mathbb{Q}_{\ell}}\ } \mathrm{Mod}(\mathbb{Q}_{\ell})
  \qquad A \rightsquigarrow T_{\ell}(A) \otimes_{\mathbb{Z}_{\ell}} \mathbb{Q}_{\ell}
\]
\[
  \mathcal{VA}^{\circ} \xrightarrow{\ H^{1}\ } \mathrm{Mod}(\mathbb{Q}_{\ell})
  \qquad A \rightsquigarrow H^{1}(A, \mathbb{Q}_{\ell}) =
  \mathrm{Hom}_{\mathbb{Q}_{\ell}}(T_{\mathbb{Q}_{\ell}}(A), \mathbb{Q}_{\ell})
\]
foncteurs \emph{fidèles}, \emph{additifs}.
\note{un astérisque, en exposant de la première source $\mathcal{VA}$, paraît effacé ; l'indice de $T$ au-dessus de la première flèche surcharge un premier indice raturé ; de
même sous le $\mathrm{Hom}$ de la seconde ligne, où un indice $\mathbb{Z}_\ell$ semble remplacé
par $\mathbb{Q}_\ell$.}
\[
  H^{*}(A) = {\textstyle\bigwedge^{*}} H^{1}(A)
\]

\begin{tikzcd}[column sep=huge]
  \mathcal{VA}^{\circ} \arrow[r, "*"] \arrow[d, "(H^{1})^{\circ}"'] & \mathcal{VA} \arrow[d, "H^{1}"] \\
  \mathrm{Mod}\,\mathbb{Q}_{\ell}^{\circ} \arrow[r, "{\mathrm{Hom}(\cdot,\, \mathbb{Q}_{\ell}(-1))}"'] & \mathrm{Mod}\,\mathbb{Q}_{\ell}
\end{tikzcd}

$A \rightsquigarrow A^{*}$ (équivalences) ; commutatif à isom. près.

\struck{$Z^{i}(A)$ \ill{}}
\[
  \mathcal{VA} \xrightarrow{\ \mathcal{C}^{*}\ } (\mathrm{Ab}) \qquad
  A \rightsquigarrow \mathcal{C}^{*}(A) \quad
  \text{anneau de Chow (des cycles mod \emph{équiv. alg})}
\]
\[
  \mathcal{C}^{i} \longrightarrow H^{2i} \otimes \mathbb{Q}_{\ell}(i) =
  {\textstyle\bigwedge^{2i}} H^{1} \otimes \mathbb{Q}_{\ell}(i) \qquad
  \mathcal{C}^{i}(A) \xrightarrow{\ \gamma_{A}^{2i}\ } H^{2i}(A, \mathbb{Q}_{\ell})(i)
\]
compatible avec structure d'anneau.
\note{le $\mathcal{C}$ ronde de $\mathcal{C}^*$ et de $\mathcal{C}^i$ surcharge chaque fois une
première lettre ($Z$ ?).}

2
\[
  \begin{array}{rcl}
  H^{2n}(A) \simeq {\textstyle\bigwedge^{2n}} H^{1}(A) & \xrightarrow[\ \sim\ ]{\eta_{A}} & \mathbb{Q}_{\ell}(-n) \\[2pt]
  H^{2n}(A^{*}) \simeq {\textstyle\bigwedge^{2n}} H^{1}(A^{*}) & \xrightarrow[\ \sim\ ]{\eta_{A^{*}}} & \mathbb{Q}_{\ell}(-n) \\[2pt]
  \wr & & \wr\ \text{(can.)} \\
  \mathbb{Q}_{\ell}(2n) & & \mathbb{Q}_{\ell}(-2n)
  \end{array}
\]
\emph{isom canoniques} \quad id ? \quad vérifier s'il n'y a pas un signe !
\note{sous la seconde ligne, deux isomorphismes verticaux mènent l'un à « $\mathbb{Q}_\ell(2n)$ »
(lecture incertaine, surchargée), l'autre, marqué « can. » dans un cercle, à
$\mathbb{Q}_\ell(-2n)$ ; des croix au-dessus de $H^1(A^*)$ et du second $\mathbb{Q}_\ell(-n)$
semblent renvoyer l'un à l'autre. « id » est suivi d'un signe en équerre.}

\page{79}
\note{la page 78, entre 77 et 79, est une page dactylographiée étrangère à ces notes ;
l'alinéa 3) ci-dessous suit les alinéas (1) et 2 de la page 77.}

3) Soit $u : A' \to A$ un \struck{isogénie} homom., avec $\dim A = \dim A' = n$, alors
\[
  \begin{array}{ccc}
  u^{2n} : H^{2n}(A) & \longrightarrow & H^{2n}(A') \\
  \wr & & \wr \\
  \mathbb{Q}_{\ell}(-n) & & \mathbb{Q}_{\ell}(-n)
  \end{array}
\]
est donné par la multiplication par $\deg u = d$ (nul ssi $u(A') \neq A$, i.e. $u$ n'est pas une
\emph{isogénie}). En d'autres termes, définissant \struck{\ill{}}
\[
  u_{*} : H^{*}(A) \longrightarrow H^{*}(A')
\]
par transposition (\uncertain{attention}, il y a que $A$ et $A'$ sont de même dim) on trouve
\[
  u_{*}(1_{A}) = d \cdot 1_{A'}
\]
Supposons que \struck{les} \struck{degrés} dim $n$, $n'$ de $A$, $A'$ soient quelconques, on
définit par transposition
\[
  H^{i}(A) \longrightarrow H^{i+2(n'-n)}(A')(n'-n)
\]
\struck{et on \ill{}} qui se déduit directement de la connaissance de
\[
  u^{1} : H^{1}(A') \longrightarrow H^{1}(A) \qquad
  (\text{d'où } u^{*} : H^{*}(A') \to H^{*}(A) \text{ par } {\textstyle\bigwedge^{*}})
\]
et des isom.\ $\eta_{A}$ et $\eta_{A'}$, par transposition. Notons que

\page{80}
\struck{\ill{}}
\[
  u_{*}(1_{A}) \in H^{2(n'-n)}(A')(n'-n)
\]
est la classe du cycle image $u(A) \in \mathcal{C}^{n'-n}(A')$ (de façon générale, $u_{*}$
compatible avec \uncertain{les} \uncertain{images} \uncertain{directes} sur les cycles …), donc
est nulle si $u(A)$ est de dim $< n$, et si $u(A)$ est de dim $n$, et si $B$ désigne la VA
\struck{\ill{}} image (au sens \uncertain{faisceautique}) et $d$ le degré de
\uncertain{projection} $A \to B$, alors $u_{*}(1_{A})$ est égal à :
$d \cdot \gamma^{2(n'-n)}_{A'}(B)$ (classe de cycle associée).
\note{« classe de cycle associée » est écrit sous $\gamma^{2(n'-n)}_{A'}(B)$ et rattaché par une
flèche. Le texte porte $u(A) \in \mathcal{C}^{n'-n}(A)$ ; la cohérence demande $(A')$.}

En termes de $H^{1}(A'^{*}) \to H^{1}(A)$ : on regarde la transposée \ill{} des
$\mathbb{Q}_{\ell}$ \struck{\ill{}},
\[
  u_{1} : T_{\mathbb{Q}_{\ell}}(A) \longrightarrow T_{\mathbb{Q}_{\ell}}(A'),
\]
\struck{on prend} s'il n'est pas injectif on trouve $0$ ; s'il est injectif on trouve un
multivecteur \struck{total} de degré $2(n'-n) = \dim \operatorname{Coker} u_{1}$, qui définit
une forme $2(n'-n)$-fois linéaire \uncertain{alternée} sur $T_{\mathbb{Q}_{\ell}}(A')$
\note{la phrase se poursuit sur la page suivante (lot suivant).}

\end{document}
